% NoVectDB / CookiX.
%
% Every quantitative claim in this paper is produced by a command in this repository
% and the command is named next to the number. See docs/paper/CLAIMS.md for the claim
% discipline this source follows.
%
% Build:  pdflatex novectdb && pdflatex novectdb

\documentclass[11pt]{article}

\usepackage[utf8]{inputenc}
\usepackage[T1]{fontenc}
\usepackage{amsmath,amssymb,amsthm}
\usepackage{booktabs}
\usepackage{algorithm}
\usepackage{graphicx}
\usepackage{hyperref}
\usepackage[margin=1in]{geometry}

\newtheorem{theorem}{Theorem}[section]
\newtheorem{proposition}[theorem]{Proposition}
\newtheorem{corollary}[theorem]{Corollary}
\newtheorem{lemma}[theorem]{Lemma}
\theoremstyle{definition}
\newtheorem{definition}[theorem]{Definition}
\newtheorem{remark}[theorem]{Remark}

\newcommand{\cmd}[1]{\texttt{\small #1}}

\title{NoVectDB: Typed Relational Retrieval as a Complement to Vector Search\\
\large CookiX: A Reference Implementation, and What It Does and Does Not Buy}

\author{Ahmed Hafdi\\Independent Research, CookiX Project}

\date{August 2026}

\begin{document}
\maketitle

\begin{abstract}
Retrieval-augmented generation grounds language models in external knowledge
predominantly through vector similarity. That works well for topical lookup but
degrades on \emph{relational} questions --- ones whose answer depends on a typed,
directed connection (\emph{prevents}, \emph{causes}, \emph{depends on}) rather than on
proximity --- and it returns a scalar distance in place of a justification. We
describe \textbf{NoVectDB}, a design position rather than a replacement: retrieval
substrates should carry typed directed edges alongside vectors, and should return the
\emph{path} that justifies an answer. We define the Knowledge Object
$K=(V,E,\mathcal{T},\mathcal{S})$, give a composite distance over the resulting graph
in which every term is ablatable, and present \textbf{CookiX}, an open-source
single-node reference implementation in pure Python.

We report what we measured, including where it loses. On 2WikiMultiHopQA
\cite{ho2020constructing}, typed multi-hop traversal reaches hits@10 $0.580$ against
Okapi BM25's \cite{robertson2009probabilistic} $0.386$ ($+50\%$ relative)
\emph{under oracle entity linking}, and recovers the gold relation chain on most
answered questions. Without the oracle, a surface-form linker achieves $59.5\%$ link
accuracy and end-to-end hits@10 $0.378$ --- \emph{parity with, not victory over},
BM25's $0.386$. Entity linking and triple extraction, not traversal, are the binding
constraints: measured per-edge extraction recall is $0.444$, which compounds to
$\approx4\%$ survival for a four-hop chain.

Neither exploratory layer earns its place on our evidence. The topological signature
($\mathcal{T}$) changes no retrieval metric we measure. For the sheaf layer
($\mathcal{S}$) we introduce a discriminative test --- does the composition residual
rank a gold evidence chain below a corrupted one on held-out data? --- and report a
\textbf{negative result} across nine configurations: learned orthogonal restriction
maps do not separate chains better than an identity-map control by more than sampling
noise. We identify the mechanism: the residual is blind to which of several genuine
paths connects two endpoints, and that is exactly the discrimination retrieval
re-ranking requires. We also state explicitly that no dense vector database has yet
been benchmarked against CookiX, and that doing so is the most important outstanding
experiment.
\end{abstract}

\noindent\textbf{Keywords:} knowledge graphs, multi-hop retrieval, retrieval-augmented
generation, explainable retrieval, sheaf theory, topological data analysis, negative
results

\section{Introduction}

The dominant pattern for grounding language models in external knowledge is
retrieval-augmented generation \cite{lewis2020retrieval}, which relies largely on
embedding text into $\mathbb{R}^n$ and retrieving neighbours by cosine or Euclidean
distance. This is effective for topical retrieval. It is weaker for a specific class of
question, and it is that class we address.

\subsection{Motivating example: the umbrella problem}

Consider five concepts: \emph{raincoat}, \emph{rain}, \emph{water}, \emph{umbrella},
\emph{storm}. A sentence transformer maps them to points that cluster tightly. The
query ``what prevents rain from reaching the coat?'' retrieves all five with similar
scores. The correct answer, \emph{umbrella}, depends on a directed edge
\[
  \textit{umbrella} \xrightarrow{\ \textit{prevents}\ } \textit{rain}
  \xrightarrow{\ \textit{causes}\ } \textit{wet\_coat},
\]
and no ranking by a \emph{symmetric distance} can encode the verb ``prevents''
(Proposition~\ref{prop:gap}).

The scope of that statement matters. It is a limitation of \emph{distance-ranked}
retrieval, not of vector representations as such. Knowledge-graph embedding methods
with \emph{asymmetric scoring functions} --- TransE \cite{bordes2013translating},
ComplEx \cite{trouillon2016complex}, RotatE \cite{sun2019rotate} --- were designed
precisely to represent directed relations in vector space, and nothing here refutes
them. The gap we exploit is between typed traversal and \emph{cosine-similarity
retrieval over text embeddings}, which is what production RAG stacks deploy.

\subsection{Contributions}

\begin{enumerate}
  \item We formalise the NoVectDB data model (Section~\ref{sec:paradigm}): Knowledge
  Objects, the graph substrate, and a composite distance whose every term is ablatable.
  \item We give the retrieval pipeline and its complexity
  (Section~\ref{sec:cookix}), as implemented.
  \item We report reproducible benchmarks (Section~\ref{sec:experiments}) on a
  synthetic corpus and on 2WikiMultiHopQA, separating the \emph{engine under oracle
  linking} from the \emph{end-to-end system}, and reporting the latter as parity rather
  than a win.
  \item We quantify extraction as the binding constraint on multi-hop accuracy
  (Section~\ref{sec:extraction}), including the compounding model.
  \item We report a \textbf{negative result} for the sheaf layer
  (Section~\ref{sec:sheaf-negative}), with a structural diagnosis of why it fails at
  the task retrieval needs.
\end{enumerate}

\section{Where distance-ranked retrieval struggles}

\subsection{Concentration of measure, and how much it actually matters}

\begin{theorem}[Distance concentration \cite{vershynin2018high}]
Let $X_1,\dots,X_m$ be i.i.d.\ random vectors drawn uniformly from the unit ball
$B^n(0,1)\subset\mathbb{R}^n$. Then for $i\neq j$,
\[
  \Pr\!\left[\left|\,\|X_i-X_j\| - \sqrt{2}\,\right| > \varepsilon\right]
  \;\le\; 2\exp\!\left(-c\,n\varepsilon^4\right),
\]
for an absolute constant $c>0$.
\end{theorem}

\begin{remark}[Scope]
\label{rem:concentration}
This should not be read as a practical indictment of vector retrieval. The theorem
concerns vectors drawn \emph{uniformly from a ball}; learned text embeddings are
strongly anisotropic and concentrate near low-dimensional manifolds, which is why
cosine retrieval remains discriminative at $n=768$--$3072$ in practice. Concentration
of measure is an asymptotic caution about isotropic data, not the failure mode
practitioners observe. Our argument rests instead on
Section~\ref{sec:semantic-gap}, which is independent of it.
\end{remark}

\subsection{The semantic gap: a symmetric distance cannot encode an asymmetric relation}
\label{sec:semantic-gap}

\begin{definition}[Semantic gap]
Let $\phi:\mathcal{D}\to\mathbb{R}^n$ be an embedding of a corpus $\mathcal{D}$, and
let $\rho$ be a target relation (e.g.\ \emph{causes}, \emph{contradicts}). The
embedding exhibits a \emph{semantic gap} with respect to $\rho$ if there exist
$a,b,c\in\mathcal{D}$ with
\[
  \rho(a,b)=\text{true},\quad \rho(a,c)=\text{false},\quad
  \|\phi(a)-\phi(b)\| \;\ge\; \|\phi(a)-\phi(c)\|.
\]
\end{definition}

\begin{proposition}
\label{prop:gap}
Let $d$ be any metric on $\mathbb{R}^n$ and let retrieval rank candidates by
$d(\phi(a),\cdot)$. For any non-symmetric relation $\rho$ there exist corpora on which
this ranking exhibits a semantic gap.
\end{proposition}

\begin{proof}
A non-symmetric $\rho$ distinguishes the ordered pairs $(a,b)$ and $(b,a)$, whereas
any metric satisfies $d(x,y)=d(y,x)$. A ranking that is a function of $d$ alone
therefore assigns the same score to both orderings and cannot separate them.
\end{proof}

\begin{remark}[What this does and does not cover]
Proposition~\ref{prop:gap} is a statement about \emph{metric-ranked} retrieval. An
asymmetric scoring function over embeddings is not a metric and escapes the argument
entirely; the knowledge-graph embedding literature
\cite{bordes2013translating,trouillon2016complex,sun2019rotate} is built on exactly
that escape. The claim we make is narrower and, we think, still worth making: the
deployed default --- rank text embeddings by cosine --- cannot represent direction.
\end{remark}

\subsection{Opacity}

Vector retrieval returns a scalar. There is no path and no typed connection, so when a
RAG system answers incorrectly, attributing the error to a retrieval decision requires
inspecting inner products. Returning the justifying path is the capability we consider
most defensible, and Section~\ref{sec:experiments} reports it as \texttt{path\_match}.

\section{The NoVectDB data model}
\label{sec:paradigm}

NoVectDB stands for \emph{Not Only Vector Database}, after NoSQL: vectors are one tool
among several.

\begin{definition}[Knowledge Object]
A Knowledge Object is a quadruple $K=(V,E,\mathcal{T},\mathcal{S})$:
\begin{itemize}
  \item $V\in\mathbb{R}^n$ (optional): an embedding or content signature.
  \item $E=\{(r_i,K_i,w_i)\}$: typed, directed, weighted edges, with $r_i$ drawn from
  an extensible vocabulary $\mathcal{R}=\{\textit{causes}, \textit{is\_a},
  \textit{part\_of}, \textit{prevents}, \textit{contradicts},\dots\}$ carrying
  algebraic properties (symmetry, transitivity, inverse).
  \item $\mathcal{T}\in\mathbb{R}^m$ (optional, \emph{no measured benefit}): a
  topological signature from persistent homology of the local neighbourhood.
  \item $\mathcal{S}$ (optional, \emph{negative result}): a sheaf stalk with linear
  restriction maps per edge.
\end{itemize}
$V$ is optional: pure topological--relational storage is a valid configuration.
\end{definition}

\begin{definition}[Composite distance]
\label{def:distance}
For Knowledge Objects $K_a,K_b$,
\begin{equation}
  d(K_a,K_b) \;=\; \alpha\cdot d_{\mathrm{geo}}(a,b)
  \;+\; \beta\cdot\bigl(1-\mathrm{TVS}(\mathcal{T}_a,\mathcal{T}_b)\bigr)
  \;+\; \gamma\cdot\lVert \mathcal{S}_a \circ_\pi \mathcal{S}_b \rVert,
  \label{eq:distance}
\end{equation}
with $\alpha+\beta+\gamma=1$, where $d_{\mathrm{geo}}$ is the minimum-weight path
distance, $\mathrm{TVS}$ compares persistence signatures, and the third term is the
sheaf composition residual along path $\pi$.
\end{definition}

Setting $\beta=\gamma=0$ recovers pure typed traversal. Every result in
Section~\ref{sec:experiments} is reported per ablation, and on our measurements
$\beta,\gamma>0$ never helps. We regard that as the point of building the distance this
way: the exotic terms are hypotheses held where they can be switched off and measured.

\section{Exploratory layer I: topological signatures}

\begin{definition}[Topological signature]
For a Knowledge Object $K$ with $r$-hop neighbourhood $N_r(K)$, build the
Vietoris--Rips complex on the shortest-path distance matrix using edge weights as
filtration values, and set
$\mathcal{T}(K)=\mathrm{Vectorise}\bigl(\mathrm{Barcode}(\mathrm{VR}(N_r(K)))\bigr)$.
\end{definition}

\begin{definition}[Topological vector similarity]
$\mathrm{TVS}(\mathcal{T}_a,\mathcal{T}_b)=\exp(-\lambda\,W_2(\mathrm{Dgm}_a,
\mathrm{Dgm}_b))$ for bandwidth $\lambda>0$ and 2-Wasserstein distance $W_2$
\cite{edelsbrunner2010computational}.
\end{definition}

\begin{remark}[Status: no measured benefit]
\label{rem:topo-status}
Across every benchmark in Section~\ref{sec:experiments}, enabling $\mathcal{T}$ changes
no retrieval metric outside noise. The layer remains in the implementation because it
is ablatable and costs nothing switched off, not because we have evidence for it. The
hypothesis --- that the \emph{shape} of a concept's neighbourhood carries retrieval
signal --- is not refuted by this, but it is not supported either, and we make no claim
for it.
\end{remark}

\section{Exploratory layer II: sheaf composition}
\label{sec:sheaf}

\begin{definition}[Cellular sheaf \cite{curry2014sheaves,hansen2019toward}]
A cellular sheaf $\mathcal{F}$ on a directed graph $G=(V,E)$ assigns a vector space
(\emph{stalk}) $\mathcal{F}(v)\cong\mathbb{R}^{d}$ to each vertex and a linear
\emph{restriction map} to each edge. In CookiX a stalk encodes an object's local
semantic frame and a typed edge's restriction map encodes how meaning transforms across
that relation.
\end{definition}

\begin{definition}[Composition residual]
For a path $\pi=(v_0,e_1,v_1,\dots,e_k,v_k)$, the composed map is
$S_\pi=\mathcal{F}_{e_k}\circ\cdots\circ\mathcal{F}_{e_1}$ and the residual is
\begin{equation}
  \mathrm{res}(\pi) \;=\; \lVert S_\pi(x_{v_0}) - x_{v_k}\rVert.
  \label{eq:residual}
\end{equation}
The intended reading is that a low residual indicates a coherent chain.
\end{definition}

\begin{remark}[Structural blindness]
\label{rem:blind}
Equation~\eqref{eq:residual} reads only the source stalk, the relation sequence, and
the target stalk. Intermediate vertices never enter it. The residual therefore
\emph{cannot} distinguish two paths that share endpoints and relation sequence, and as
Section~\ref{sec:sheaf-negative} shows empirically, it does not distinguish paths that
share endpoints with a \emph{different} relation sequence either. This is a property of
the definition, not of a particular fit, and it is the reason the layer fails at
re-ranking.
\end{remark}

\section{Retrieval-theoretic analysis}
\label{sec:theory}

\begin{definition}[Relational query]
A relational query $q=(K_{\mathrm{start}},\rho_1,\dots,\rho_k)$ seeks objects reachable
from $K_{\mathrm{start}}$ along a chain of typed relations.
\end{definition}

\begin{proposition}[Conditional advantage]
\label{prop:conditional}
Let $\mathrm{Prec}_V(q)$ be the precision@$k$ of metric-ranked vector retrieval and
$\mathrm{Prec}_N(q)$ that of typed traversal, for a query of hop count $h$. Suppose
each required edge is present and correctly typed in the graph independently with
probability $p$, and the head entity is correctly identified with probability $\ell$.
Then
\begin{equation}
  \mathrm{Prec}_N(q) \;\lesssim\; \ell \cdot p^{\,h},
  \label{eq:compounding}
\end{equation}
and traversal outperforms metric-ranked retrieval only when
$\ell\,p^{\,h} > \mathrm{Prec}_V(q)$.
\end{proposition}

\begin{proof}[Sketch]
Traversal recovers the target only if the anchor is correct and every hop's edge is
present and correctly typed; independence gives the product. Correct traversal is
precision~$1.0$ for single-hop typed lookup, so the bound is tight at $h=1$,
$\ell=p=1$.
\end{proof}

\begin{remark}[The condition is not satisfied in the open-domain setting]
\label{rem:condition-fails}
Equation~\eqref{eq:compounding} is the honest centre of this paper, and it cuts against
us. At our measured values --- $p\approx0.444$ from Section~\ref{sec:extraction} and
$\ell\approx0.595$ from Section~\ref{sec:endtoend} --- two-hop survival is
$0.595\times0.444^2\approx0.117$, and the inequality fails against a BM25 baseline at
$0.386$. This is not a discrepancy to be explained away: it is exactly what the
end-to-end measurement shows, where CookiX reaches parity rather than beating BM25.
Under oracle linking ($\ell=1$) with a gold-triple graph ($p=1$) the same inequality
predicts the advantage we do observe. The advantage is real and the condition for it is
specific.
\end{remark}

\begin{corollary}[Dimension independence]
Because $d_{\mathrm{geo}}$ is defined on a discrete graph and $\mathrm{TVS}$ derives
from topological invariants, neither depends on an ambient embedding dimension.
\end{corollary}

\begin{remark}
This is a categorical difference, not a demonstrated benefit; given
Remark~\ref{rem:concentration} it should not be read as an advantage over deployed
vector retrieval.
\end{remark}

\section{CookiX: architecture}
\label{sec:cookix}

CookiX is a document-oriented embedded library, in pure Python (3.10+) with NumPy and
NetworkX as its only core dependencies.

\begin{itemize}
  \item \textbf{Ingestor.} Accepts structured documents, or free text via a swappable
  extractor: a deterministic rule-based keyword extractor, or an LLM-backed extractor
  through an API. Extraction quality is the ceiling on everything downstream
  (Section~\ref{sec:extraction}), which is why it is a first-class component.

  \item \textbf{Storage.} Three interchangeable backends behind one interface, held to a
  shared parity test battery: an in-memory NetworkX store; a crash-safe durable store in
  pure Python (write-ahead log with per-record CRC for torn-write tolerance, atomic
  snapshots via temp-file rename, single-\texttt{fsync} transactions, re-entrant write
  lock); and an embedded Kùzu property graph.

  \item \textbf{TopoIndex.} Approximate nearest-neighbour search over persistence
  signatures by cosine LSH, with an exact fallback and a built-in recall measure.

  \item \textbf{Query engine.} The pipeline of Algorithm~\ref{alg:pipeline}. Geodesic
  search is a settle-once Dijkstra --- a node is never re-expanded once its minimum-cost
  path is fixed --- with early exit when a specific target is requested.

  \item \textbf{Interfaces.} An embedded Python API; an HTTP server with API-key auth,
  role separation, rate limiting, resource caps, read-only mode, health probes and
  Prometheus metrics; and a dependency-free typed client over a versioned wire API.
\end{itemize}

\begin{algorithm}
\caption{NoVectDB retrieval pipeline}
\label{alg:pipeline}
\begin{enumerate}
  \item $I \leftarrow \textsc{IntentParse}(q)$ \hfill (anchor / relation / target)
  \item $S_0 \leftarrow \textsc{DeterministicLookup}(G, I)$ \hfill (exact typed-edge
  match; inverse relations resolved virtually)
  \item \textbf{if} a bare relation query \textbf{then return} $S_0[{:}k]$
  \item $S_1 \leftarrow \textsc{GeodesicSearch}(G, I.\mathrm{anchor}, h_{\max})$
  \hfill (settle-once Dijkstra, early exit on target)
  \item $S_2 \leftarrow \textsc{Rank}(S_1, d)$ \hfill (Eq.~\eqref{eq:distance}, terms
  enabled per ablation mode)
  \item \textbf{return} $S_2[{:}k]$
\end{enumerate}
\end{algorithm}

\begin{proposition}[Query complexity]
The pipeline runs in $O(|V_{\mathrm{local}}| + |E_{\mathrm{local}}|\log
|E_{\mathrm{local}}|)$, where the local subscripts denote the explored frontier ---
bounded by degree and hop limit, not by total object count.
\end{proposition}

\section{Experiments}
\label{sec:experiments}

Every number in this section is produced by a named command in this repository, with a
fixed seed. Ablation modes are reported separately throughout.

\subsection{Synthetic relational corpus}

\cmd{cookix eval --seed 0 --worlds 40 --k 5}. The corpus is built so that each entity's
text describes the entity itself, never its relations, and entities within a world share
a topical adjective --- so a content retriever genuinely retrieves the right
neighbourhood, while the relational answer lives only in the typed edges. 240 documents,
160 queries (single-hop forward, single-hop inverse, multi-hop, contradiction).

\begin{table}[h]
\centering
\begin{tabular}{lrrrrrr}
\toprule
retriever & hits@1 & hits@5 & P@5 & R@5 & MRR & path\_acc \\
\midrule
random           & 0.006 & 0.019 & 0.004 & 0.019 & 0.010 & 0.000 \\
lexical-tfidf    & 0.250 & 0.750 & 0.150 & 0.750 & 0.375 & 0.000 \\
cookix-graph     & 1.000 & 1.000 & 0.200 & 1.000 & 1.000 & 1.000 \\
cookix-topo      & 1.000 & 1.000 & 0.200 & 1.000 & 1.000 & 1.000 \\
cookix-sheaf     & 1.000 & 1.000 & 0.200 & 1.000 & 1.000 & 1.000 \\
cookix-reasoning & 1.000 & 1.000 & 0.200 & 1.000 & 1.000 & 1.000 \\
\bottomrule
\end{tabular}
\caption{Synthetic corpus, seed 0. CookiX saturates, so the exotic layers have no
headroom to demonstrate anything and we claim nothing from their tying.}
\end{table}

\begin{remark}[What this table is worth]
This is a corpus we designed, and saturation at $1.000$ is evidence that the
implementation does what it is built to do --- not evidence about the world. The
informative part is the baseline: the lexical retriever reaches $R@5=0.75$ while scoring
$\text{hits@}1=0.25$ and $\text{path\_acc}=0$. Content similarity finds the right
neighbourhood, cannot pick the relationally correct entity out of it, and structurally
cannot produce a path.
\end{remark}

\subsection{2WikiMultiHopQA under oracle entity linking}

\cmd{cookix eval --dataset 2wiki --path dev.json --k 10}. One global knowledge graph is
built from every example's gold evidence triples, so traversal faces real distractor
edges from thousands of other questions. Baseline is Okapi BM25 over the same
paragraphs. First 2{,}000 dev examples, 1{,}802 evaluable.

\begin{table}[h]
\centering
\begin{tabular}{lrrr}
\toprule
retriever & hits@10 & MRR & path\_match \\
\midrule
BM25 (strong lexical) & 0.386 & 0.239 & n/a \\
cookix-graph          & \textbf{0.580} & 0.282 & 0.579 \\
cookix-reasoning      & 0.580 & 0.283 & 0.579 \\
\bottomrule
\end{tabular}
\caption{Oracle entity linking: the question's head entity is given, so this measures
the reasoning engine in isolation. \texttt{path\_match} is unavailable for BM25 by
construction --- it returns passages, not paths.}
\end{table}

This is the strongest evidence for the typed-traversal core, in a setting standard in
KG-QA. It is \emph{not} an end-to-end open-domain number, and we do not present it as
one. Note also that \texttt{reasoning} ties \texttt{graph}: the re-ranking layers
reorder candidates without changing which answers are reachable.

\subsection{Dropping the oracle: the end-to-end result}
\label{sec:endtoend}

\cmd{cookix eval --dataset 2wiki --no-oracle --linker \{surface,bm25\}}.

\begin{table}[h]
\centering
\begin{tabular}{lrrr}
\toprule
anchor source & link accuracy & CookiX hits@10 & MRR \\
\midrule
oracle (given)              & 100\% & 0.580 & 0.282 \\
surface linker              & \textbf{59.5\%} & \textbf{0.378} & 0.183 \\
BM25 linker (context match) & 50.1\% & 0.340 & 0.168 \\
\midrule
BM25 \emph{retriever} (no linking needed) & --- & 0.386 & 0.239 \\
\bottomrule
\end{tabular}
\caption{End-to-end, same 2{,}000 examples. A correct-signal linker lifts CookiX from
below BM25 to parity. It is not a win.}
\end{table}

Stated without spin: the engine dominates when linked correctly; matching entity names
against the question beats matching paragraph content ($50.1\%\to59.5\%$ link accuracy,
$0.340\to0.378$ hits@10); and that lifts CookiX from \emph{losing} to \emph{parity}
($0.378$ vs $0.386$) while still returning a reasoning path. Roughly $40\%$ of questions
still begin from the wrong anchor and cannot recover. Clearing $\approx70\%$ link
accuracy is what would flip the comparison. An LLM-assisted linker is implemented and
unit-tested; we report no number for it because we have not run the keyed benchmark.

\subsection{Extraction is the binding constraint}
\label{sec:extraction}

\cmd{cookix eval --extraction}. Rule-based extractor against 16 hand-annotated sentences
(18 gold triples):

\begin{table}[h]
\centering
\begin{tabular}{lrrrr}
\toprule
extractor & precision & recall & F1 & relation accuracy \\
\midrule
rule-based & 0.615 & 0.444 & 0.516 & 1.000 \\
\bottomrule
\end{tabular}
\caption{Relation \emph{typing} is perfectly reliable --- whenever the extractor spans
the right entities it labels the relation correctly. Free-text \emph{coverage} is the
bottleneck.}
\end{table}

Projecting per-edge recall through Eq.~\eqref{eq:compounding} with $\ell=1$:

\begin{table}[h]
\centering
\begin{tabular}{lrrrr}
\toprule
& 1-hop & 2-hop & 3-hop & 4-hop \\
\midrule
rule-based & 0.444 & 0.198 & 0.088 & 0.039 \\
\bottomrule
\end{tabular}
\caption{A four-hop chain survives about $4\%$ of the time.}
\end{table}

The misses are the realistic cases: out-of-vocabulary synonyms (``leads to'', ``ward
off''), sentences carrying two relations where a keyword splitter emits only the first,
and preposition boundaries. How much of this ceiling an LLM extractor buys back is open;
we have not measured it.

\subsection{A negative result for the sheaf layer}
\label{sec:sheaf-negative}

\cmd{cookix eval --sheaf-probe --path dev.json --dim 32}.

An absolute residual is not a meaningful test of this layer. Fitting restriction maps on
data generated to be sheaf-consistent and observing that the residual drops establishes
only that orthogonal Procrustes recovers a linear map when a linear map produced the
data. The question that bears on retrieval is discriminative.

\paragraph{Design.} On held-out 2WikiMultiHopQA chains, we ask whether the residual
ranks a gold evidence chain below a corrupted one. Stalks are derived from entity
\emph{text} (TF--IDF, randomly projected, L2-normalised) rather than from an entity-id
hash, since a hash-seeded stalk carries no semantics and no result computed against it
would mean anything. Corruptions vary only what Eq.~\eqref{eq:residual} can see: a
swapped relation label, a swapped source or target, a shuffled relation order, and a
\emph{detour} --- a genuine alternative path between the same endpoints. Three map
families are scored:

\begin{itemize}
  \item \emph{placeholder} --- the shipped random orthogonal maps. The null hypothesis.
  \item \emph{identity} --- every map is $I$, collapsing the residual to endpoint
  distance $\lVert x_a-x_b\rVert$ and ignoring the relation chain entirely. This is the
  decisive control: it asks whether the sheaf machinery does anything that embedding
  distance does not.
  \item \emph{learned} --- orthogonal Procrustes fitted on training-split edges only,
  with every test-chain triple withheld.
\end{itemize}

The decision rule was fixed before running: \emph{learned} must beat \emph{both}
controls by more than the sampling-noise floor $1/\sqrt{n}$. Beating the placeholder
alone would show only that fitting beats not fitting.

\paragraph{Result.} 12{,}576 dev examples yield 6{,}459 connected chains; 1{,}938 are
held out and scored, with 26{,}126 training edges over 34 relations (median 92
edges/relation against $\dim=32$, so the fit is not underdetermined). Noise floor
$0.023$.

\begin{table}[h]
\centering
\begin{tabular}{lrrrrrr}
\toprule
maps & pooled AUC & detour & rel.\ swap & shuffle & src.\ swap & tail swap \\
\midrule
placeholder        & 0.509 & 0.498 & 0.504 & 0.524 & 0.502 & 0.508 \\
identity (control) & 0.530 & 0.463 & 0.500 & 0.523 & 0.563 & 0.536 \\
learned            & 0.552 & 0.481 & 0.550 & 0.530 & 0.565 & 0.561 \\
\bottomrule
\end{tabular}
\caption{AUC $=\Pr(\text{gold residual} < \text{corrupted residual})$; $0.5$ is chance.
\emph{learned} $-$ \emph{identity} $=0.022$, below the $0.023$ noise floor.}
\end{table}

Across nine configurations ($\dim\in\{8,16,32\}$, three seeds) the
\emph{learned}\,--\,\emph{identity} gap ranges from $+0.004$ to $+0.018$ and is
\textbf{below the noise floor in all nine}. We therefore report a negative result: on
real multi-hop data, learned linear restriction maps over lexical stalks do not
discriminate coherent from corrupted reasoning chains better than endpoint distance.

\paragraph{Diagnosis.} The \emph{detour} column is $0.481$ --- chance. The residual
cannot tell a gold chain from another genuine path between the same endpoints. This
explains mechanically why $\gamma>0$ never improves retrieval in our benchmarks:
retrieval candidates \emph{are} all real paths in the graph, i.e.\ all detours, which is
precisely the discrimination the residual lacks (cf.\ Remark~\ref{rem:blind}). The layer
was moderately effective in a controlled setting at detecting wrong relation labels and
wrong endpoints --- a \emph{verification} question, not a retrieval one.

\paragraph{What remains open.} Stalks here are lexical projections, not neural sentence
embeddings, and the maps are linear and closed-form rather than jointly learned as in
neural sheaf diffusion \cite{bodnar2022neural}. Those are the substantive untried
variables. We claim only what we tested: the linear-maps-over-lexical-stalks formulation
does not work, and the burden of proof now sits with the neural variant.

\subsection{Performance and scale}

\cmd{cookix eval --perf}, 240 documents, 160 queries $\times$ 3 repeats,
single-threaded: median $0.06$\,ms graph-only, $0.07$\,ms with all layers
($\approx13{,}000$ vs $\approx6{,}500$ q/s). \cmd{cookix eval --scale}, random typed
graph, out-degree 4, four-hop traversals:

\begin{table}[h]
\centering
\begin{tabular}{lrrrr}
\toprule
objects & ingest & peak memory & median query & p95 \\
\midrule
1{,}000  & 0.08\,s & 3\,MB   & 1.5\,ms & 2.8\,ms \\
10{,}000 & 0.9\,s  & 30\,MB  & 2.0\,ms & 3.1\,ms \\
50{,}000 & 6.4\,s  & 155\,MB & 2.2\,ms & 6.1\,ms \\
\bottomrule
\end{tabular}
\caption{Latency stays near-flat across a $50\times$ increase in graph size, because
traversal cost is bounded by the local reachable frontier rather than total object
count. Single-threaded pure Python on one machine; the largest graph measured is
$5\times10^4$ objects.}
\end{table}

\cmd{cookix loadtest}: 8 concurrent clients against the real HTTP server over a
10{,}000-object graph sustained $\approx130$ req/s with $0$ errors in $2{,}613$ requests
(median $59$\,ms, p99 $102$\,ms) and non-monotonic memory. That throughput is the honest
ceiling of a single synchronous Python process; multi-worker deployment is untested.

\subsection{No dense-retrieval baseline has been run}
\label{sec:no-vector-baseline}

We state this explicitly rather than leaving it to be inferred. CookiX has been compared
against a random floor, TF--IDF cosine, and Okapi BM25. It has \textbf{not} been
compared against a dense embedding retriever, Chroma, Pinecone, Weaviate, or GraphRAG.
TF--IDF stands in for the vector-similarity family on the grounds that a dense embedder
plugs into the same interface and behaves the same way with respect to \emph{relations}
--- it retrieves by topical proximity, not traversal --- but that is an argument, not a
measurement. Running a dense baseline is the most important outstanding experiment in
this line of work.

\section{Discussion}

\paragraph{When to use what.} For fuzzy semantic search over unstructured text, use a
vector database. For a production graph database at scale, use Neo4j or Kùzu. For
LLM-grounded summarisation over a corpus graph, use GraphRAG \cite{edge2024local}. The
niche we can defend is narrow: \emph{embedded, explainable, typed multi-hop retrieval
where the justifying path is part of the answer}, over data that is already relational
or can be made so reliably.

\paragraph{Multiplicative fragility.} Equation~\eqref{eq:compounding} deserves emphasis
as a design property, not just a measurement. Traversal accuracy is multiplicative in
extraction and linking accuracy, whereas lexical and dense retrieval degrade gracefully:
partial signal still surfaces something useful. Traversal from a wrong anchor fails
\emph{confidently}, emitting a clean typed path --- so the interpretability benefit
inverts into a liability when the graph or the anchor is wrong. A returned path should be
treated as a claim to be checked, not as proof. We think this is the most important
practical caveat in the paper.

\paragraph{Where the sheaf idea might still live.} Section~\ref{sec:sheaf-negative}
shows the residual fails at path \emph{discrimination} while retaining some ability to
detect wrong relation labels and wrong endpoints. That is a verification question --- given
a claimed triple or chain, is it consistent with the graph? --- and it requires neither
extraction nor entity linking. It is where we would look next, with neural stalks, and
without assuming the answer.

\paragraph{Limitations.} Single-node; no clustering, sharding or replication. Pure
Python; a Rust hot-path core is unstarted. No dense-retrieval baseline. Evaluated on one
external dataset; HotpotQA and MuSiQue loaders are unwritten. The LLM extractor and LLM
linker are implemented but unmeasured. Both exploratory layers are unsupported by
evidence, one of them negatively so.

\section{Conclusion}

Knowledge has direction, and metric-ranked retrieval cannot represent direction. That
observation is sound, and typed traversal exploits it: under oracle entity linking,
CookiX reaches hits@10 $0.580$ against BM25's $0.386$ on 2WikiMultiHopQA and recovers the
gold relation chain on most answered questions. But an end-to-end system must find its
own anchor and build its own graph, and there CookiX reaches parity with a lexical
baseline from 1994 rather than beating it. Extraction and linking, not traversal, are
what stand between those two numbers.

Of the two exotic layers, one has never moved a metric and the other now carries a
measured negative result with a structural explanation for it. We report this because a
hypothesis that survives only by not being tested is not worth holding. What remains is
narrower and, we think, more useful: typed relational retrieval is a real complement to
vector search for a specific class of question, it can return the path that justifies an
answer, and the honest frontier is upstream --- in turning text into edges, and questions
into anchors.

\begin{thebibliography}{99}

\bibitem{lewis2020retrieval}
P.~Lewis et al.
\newblock Retrieval-augmented generation for knowledge-intensive NLP tasks.
\newblock \emph{NeurIPS}, 2020.

\bibitem{vershynin2018high}
R.~Vershynin.
\newblock \emph{High-Dimensional Probability}.
\newblock Cambridge University Press, 2018.

\bibitem{edelsbrunner2010computational}
H.~Edelsbrunner and J.~L. Harer.
\newblock \emph{Computational Topology: An Introduction}.
\newblock American Mathematical Society, 2010.

\bibitem{curry2014sheaves}
J.~Curry.
\newblock Sheaves, cosheaves and applications.
\newblock arXiv:1303.3255, 2014.

\bibitem{hansen2019toward}
J.~Hansen and R.~Ghrist.
\newblock Toward a spectral theory of cellular sheaves.
\newblock \emph{Journal of Applied and Computational Topology}, 3(4):315--358, 2019.

\bibitem{bodnar2022neural}
C.~Bodnar, F.~Di~Giovanni, B.~Chamberlain, P.~Liò, and M.~Bronstein.
\newblock Neural sheaf diffusion: a topological perspective on heterophily and
oversmoothing in GNNs.
\newblock \emph{NeurIPS}, 2022.

\bibitem{edge2024local}
D.~Edge et al.
\newblock From local to global: a GraphRAG approach to query-focused summarization.
\newblock arXiv:2404.16130, 2024.

\bibitem{bordes2013translating}
A.~Bordes, N.~Usunier, A.~Garcia-Durán, J.~Weston, and O.~Yakhnenko.
\newblock Translating embeddings for modeling multi-relational data.
\newblock \emph{NeurIPS}, 2013.

\bibitem{trouillon2016complex}
T.~Trouillon, J.~Welbl, S.~Riedel, É.~Gaussier, and G.~Bouchard.
\newblock Complex embeddings for simple link prediction.
\newblock \emph{ICML}, 2016.

\bibitem{sun2019rotate}
Z.~Sun, Z.-H. Deng, J.-Y. Nie, and J.~Tang.
\newblock RotatE: knowledge graph embedding by relational rotation in complex space.
\newblock \emph{ICLR}, 2019.

\bibitem{ho2020constructing}
X.~Ho, A.-K. Duong~Nguyen, S.~Sugawara, and A.~Aizawa.
\newblock Constructing a multi-hop QA dataset for comprehensive evaluation of reasoning
steps.
\newblock \emph{COLING}, 2020.

\bibitem{robertson2009probabilistic}
S.~Robertson and H.~Zaragoza.
\newblock The probabilistic relevance framework: BM25 and beyond.
\newblock \emph{Foundations and Trends in Information Retrieval}, 3(4):333--389, 2009.

\end{thebibliography}

\end{document}
